% ssp245-review-slides.tex -- slide version of the ssp245 validation review
% Build: TEXINPUTS=/var/home/aoz/Dropbox/OZ-Labs/templates//: xelatex ssp245-review-slides.tex  (run twice)
\documentclass[aspectratio=169,11pt]{beamer}

\usetheme{ozlabs}

\renewcommand{\ozfootbrand}{OZ Labs}
\renewcommand{\ozfooturl}{ozstats.space}

\usepackage{booktabs}

\title{Validation Review: SSP2-4.5 Global Temperature-Attributable Burden Run}
\subtitle{204 countries \(\times\) 27 climate models \(\times\) 2022--2050}
\author{Aaron Osgood-Zimmerman}
\institute{OZ Labs LLC, for the World Bank climate-health project}
\date{August 18, 2026}

\begin{document}

\begin{frame}[plain,noframenumbering]
  \titlepage
\end{frame}

\begin{frame}{Bottom line}
  \begin{itemize}
    \item The first full production run (scenario SSP2-4.5) passes every
      arithmetic and coverage check.
    \item Against the published estimates of the methodology we implement
      (Burkart et al., \textit{The Lancet} 2021), our global total agrees to
      \ozgreentext{0.1\%}, and 159 of 204 countries fall inside the published
      95\% uncertainty intervals.
    \item An apparent large shortfall in hot countries turned out to be
      \ozcyantext{GBD's own upward revision between rounds}, not a pipeline
      defect.
    \item Main quantified limitations: risk curves end at
      32.1\,\textcelsius{}; a tropical cluster where daily temperature
      variability looks damped; a YLL definition that differs from GBD's.
    \item \ozgreentext{Nothing found blocks the remaining three scenarios.}
  \end{itemize}
\end{frame}

\begin{frame}{What this run is, in one slide}
  \begin{itemize}
    \item Goal: deaths and years of life lost (YLLs) attributable to
      non-optimal ambient temperature, projected under emissions scenarios.
    \item Method: Burkart et al.\ (2021), the study behind the Global Burden
      of Disease (GBD) treatment of temperature. We use the risk curves,
      optimal temperatures, and code released with that paper.
    \item Per country and cause: daily temperature \(\rightarrow\) relative
      risk vs.\ the local optimum \(\rightarrow\) population attributable
      fraction (PAF, the share of deaths attributable to non-optimal
      temperature) \(\rightarrow\) attributable deaths and YLLs.
    \item Inputs: CCKP 0.25-degree daily temperature (World Bank), IHME
      national mortality forecasts, released curves and optima.
    \item Scale: 204 countries, 27 climate models, 29 years (2022--2050),
      500 draws; 159{,}732 output files; run by Caspar, July--August 2026.
  \end{itemize}
\end{frame}

\ozsection{Data integrity}

\begin{frame}{All quality gates pass; three defects found and fixed}
  \begin{itemize}
    \item Heat + cold = non-optimal holds to floating-point precision, for
      PAFs and deaths; \(|\mathrm{PAF}| \le 1\) everywhere; all values
      finite; no all-zero combinations.
    \item Coverage complete: every country has all 27 models and all 29
      years.
    \item Three defective combinations surfaced and were recomputed
      (August 17, 2026), then re-verified:
    \begin{itemize}
      \item Colombia and Equatorial Guinea, one model, 2022: single-cause
        leftovers of an early benchmark run.
      \item Vietnam, one model, 2029: missing YLL file (interrupted write).
    \end{itemize}
    \item After the fix: all 159{,}732 combinations carry all 17 causes and
      YLLs.
  \end{itemize}
\end{frame}

\ozsection{Validation against published estimates}

\begin{frame}{Anchor 1: the nine countries published in the 2021 paper}
  \centering\small
  \begin{tabular}{lrrr}
    \toprule
    country & published 2019 [95\% UI] & ours 2022 & ratio \\
    \midrule
    China          & 499{,}605 [425{,}094--573{,}429] & 555{,}481 & 1.11 \\
    United States  & 110{,}588 [98{,}410--119{,}580]  & 103{,}964 & 0.94 \\
    Mexico         & 19{,}055 [16{,}111--22{,}083]    & 20{,}787  & 1.09 \\
    Brazil         & 17{,}286 [15{,}386--19{,}185]    & 14{,}383  & \alert{0.83} \\
    South Africa   & 8{,}815 [7{,}846--9{,}823]       & 9{,}575   & 1.09 \\
    Colombia       & 4{,}820 [3{,}669--6{,}148]       & 4{,}746   & 0.98 \\
    Chile          & 3{,}759 [3{,}244--4{,}267]       & 3{,}747   & 1.00 \\
    Guatemala      & 1{,}258 [976--1{,}637]           & 1{,}095   & 0.87 \\
    New Zealand    & 1{,}191 [996--1{,}361]           & 1{,}177   & 0.99 \\
    \bottomrule
  \end{tabular}
  \vskip0.6em
  \raggedright
  Eight of nine inside the published UIs; rank correlation 0.983. Brazil is
  low, almost entirely on the heat side. (Ours is 2022 modeled climate vs.\
  their 2019 observed weather; read as a magnitude check, not to the digit.)
\end{frame}

\begin{frame}{Anchor 2: all 204 countries vs.\ the 2021 paper's estimates}
  \begin{exampleblock}{Headline}
    Global total: ours 1.684M vs.\ published 1.686M attributable deaths per
    year (ratio 0.999). Country ranks correlate at 0.993. 159 of 204
    countries inside the published 95\% UIs.
  \end{exampleblock}
  \begin{itemize}
    \item Source: the paper's companion dataset (IHME GHDx, DOI
      10.6069/7QWW-PF74), which we first verified against the paper's own
      Table 1 (PAFs agree to \(\le 0.015\) percentage points).
    \item The 45 out-of-UI countries: 41 low (38 tropical, including 18
      small-island states), 4 high. About two thirds sit within 15\% of a UI
      bound, where comparison noise plausibly reaches.
  \end{itemize}
\end{frame}

\begin{frame}{The apparent hot-country shortfall vs.\ the current GBD round}
  \begin{itemize}
    \item Against GBD's \emph{current} (2023-round) 2022 estimates, our
      global total is 0.874 of theirs, and the biggest gaps are all hot
      countries: Niger 0.38, Sudan 0.44, Iraq 0.49, India 0.58, Nigeria
      0.57, Mali 0.58.
    \item We tested four candidate causes: temperature bias (ruled out where
      testable against the ERA5 reanalysis), unusual observed 2022 weather
      (ruled out; GBD's own series is nearly flat 2019 to 2022), mortality
      inputs (agree globally to 2\%), and risk-curve truncation
      (contributes, next section).
  \end{itemize}
\end{frame}

\begin{frame}{Resolution: it is GBD's between-round revision}
  \centering\small
  \begin{tabular}{lrrr}
    \toprule
    country & vs.\ current round & vs.\ 2021 paper & GBD's own revision \\
    \midrule
    Niger    & 0.38 & \ozgreentext{0.84} & \(\times\)2.27 \\
    Sudan    & 0.44 & \ozgreentext{1.31} & \(\times\)2.34 \\
    Kuwait   & 0.48 & \ozgreentext{1.12} & \(\times\)2.24 \\
    Iraq     & 0.49 & \ozgreentext{0.94} & \(\times\)1.87 \\
    India    & 0.58 & \ozgreentext{0.96} & \(\times\)1.67 \\
    Pakistan & 0.66 & \ozgreentext{0.90} & \(\times\)1.45 \\
    \bottomrule
  \end{tabular}
  \vskip0.6em
  \raggedright
  Compared against the estimates of the methodology we actually implement,
  the hot-country gap disappears (all rows inside the published UIs). GBD
  revised these countries upward by 1.3--2.3\(\times\) between rounds.
  Whether to track the 2021 methodology or the current round is a scoping
  decision, not a bug fix.
\end{frame}

\ozsection{Structural checks}

\begin{frame}{Trends make sense once composition is removed}
  \begin{itemize}
    \item Raw cold-attributable deaths \emph{rise} in most countries
      (1.33M in 2022 to 1.99M in 2050), which looks wrong under warming.
    \item Two composition effects explain it: population aging, and the
      mortality forecast shifting deaths toward cold-heavy causes at high
      latitudes.
    \item After age-sex standardization: heat share rises in 197 of 204
      countries; cold share falls in 140 of 204.
  \end{itemize}
  \begin{alertblock}{Reporting guidance}
    Headline trends should be age-standardized attributable
    \emph{fractions}; raw counts mix the climate signal with demography.
  \end{alertblock}
\end{frame}

\begin{frame}{Model spread and the YLL definition}
  \begin{itemize}
    \item No climate model dominates the ensemble; the warmest-tilted
      members (UKESM1-0-LL, HadGEM3-GC31-LL, KACE-1-0-G) are exactly the
      known high-sensitivity models.
    \item A few near-zero-burden countries flip sign across models (e.g.\
      Singapore, ensemble \(-96\) deaths, model range \(-171\) to
      \(+135\)); flag sign instability there, not spread.
    \item YLLs per attributable death: ours 13.9 years vs.\ GBD's 21.2. This
      is definitional: GBD uses its standard reference life table; we use
      national UN life expectancies. State it wherever YLLs are compared.
  \end{itemize}
\end{frame}

\ozsection{Quantified limitations}

\begin{frame}{Risk curves end at a daily mean of 32.1\,\textcelsius}
  \begin{itemize}
    \item Curves exist only over each zone's observed temperature range;
      hotter days get the endpoint risk. IHME's released code does the same:
      this is the published method, not our shortcut.
    \item Scale today: 7.9\% of world population lives where the assignment
      is clamped (2022), rising to 15--23\% by 2050 under SSP2-4.5. In the
      Gulf and Sahel, 25--40\% of person-days already exceed the endpoint.
    \item Extrapolating the endpoint log-linearly multiplies 2022 heat
      burden by 2.30 (Kuwait), 1.60 (Niger), 1.31 (Sudan), 1.00 (Haiti,
      no truncation): an assumption-based bound, not evidence.
  \end{itemize}
  \begin{alertblock}{Projection consequence}
    Beyond the endpoint, added warming adds no added risk, so hot-country
    projections are effectively lower bounds; report the extrapolation as a
    labeled sensitivity band.
  \end{alertblock}
\end{frame}

\begin{frame}{The tropical low cluster traces to the exposure product}
  \begin{itemize}
    \item The daily temperature input is NASA NEX-GDDP-CMIP6: climate models
      quantile-mapped to the GMFD observational dataset. Neither raw CMIP6
      nor the ERA5 reanalysis the 2021 estimates used.
    \item Where testable, its day-to-day spread runs 30--40\% below ERA5 in
      low-variance tropical countries; burden lives in that spread, so
      damped spread means damped burden. The 18 small-island states are the
      extreme case.
    \item Re-referencing to ERA5, tested on four countries, fixes Rwanda but
      moves the Philippines the \emph{wrong} way (the references disagree by
      1.9\,\textcelsius{} there). \alert{Not recommended as a production
      change on current evidence.}
    \item Takeaway: tropical country results are sensitive up to
      \(\sim\)2\(\times\) to the choice of observational reference; no
      published uncertainty interval includes that.
  \end{itemize}
\end{frame}

\begin{frame}{Other standing flags}
  \begin{itemize}
    \item \textbf{Sub-grid countries}: 29 countries span at most one
      0.25-degree grid cell; model spread \(\sim\)7\(\times\) the median,
      attributable share \(\sim\)4\(\times\) lower than resolved tropical
      peers. Flag in country tables.
    \item \textbf{Fixed optimal temperature}: the TMREL is held at its 2020
      value for all projection years (no adaptation). Defensible to 2050;
      strong for 2100; an annual derivation exists as a sensitivity.
    \item \textbf{Forecast mortality}: projected burdens mix climate with
      IHME's demographic and epidemiological forecasts; attributable
      fractions isolate the climate signal.
  \end{itemize}
\end{frame}

\ozsection{Recommendations}

\begin{frame}{Six recommendations}
  \begin{enumerate}
    \item \ozgreentext{Proceed with the remaining three scenarios}
      (SSP3-7.0, SSP5-8.5, SSP1-2.6); nothing here blocks them.
    \item State the alignment target in all outputs: these are estimates
      under the Burkart 2021 methodology. If current-round alignment is
      wanted later, obtain the current-round curves from IHME; do not tune
      ours.
    \item Report projected trends as age-standardized attributable
      fractions, counts secondary.
    \item Attach the truncation sensitivity band to hot-country projection
      results (a small local computation; no full-scale rerun).
    \item Flag the 29 sub-grid countries and the tropical low cluster in
      country-level tables.
    \item State the YLL definition wherever YLLs are shown.
  \end{enumerate}
\end{frame}

\begin{frame}{Provenance}
  \begin{itemize}
    \item Production summary tables: August 7, 2026 (re-issued August 17,
      2026, after the three-combo fix), produced by Caspar.
    \item Review analyses: August 14--17, 2026; notes in
      \texttt{ssp245-review-findings.org}, scripts and outputs in
      \texttt{output/review-ssp245/} in the project repository.
    \item External comparators: Burkart et al.\ (2021) Table 1; IHME GHDx
      dataset DOI 10.6069/7QWW-PF74 (the paper's companion estimates); GBD
      Results Tool exports (current round), July 31, 2026.
    \item A six-page written version of this review exists alongside this
      deck (\texttt{ssp245-review-summary.pdf}).
  \end{itemize}
\end{frame}

\end{document}
